\documentclass{article}
\usepackage[T1]{fontenc}
\usepackage{lmodern}
\usepackage{graphicx}
\usepackage{amsmath,amssymb}
\usepackage[a4paper,margin=2cm]{geometry}
\usepackage[absolute,overlay]{textpos}
\setlength{\TPHorizModule}{1mm}
\setlength{\TPVertModule}{1mm}
\usepackage[indonesian]{babel}
\usepackage{hyperref}
\setlength{\emergencystretch}{2em}
% unit: Halaman 19

\begin{document}

% === Halaman 19 (llm) ===

\textbf{Contoh 1:} Misalkan $a = 00001101 = x^3 + x^2 + 1$ dan $b = 00000110 = x^2 + x$\
$a$ dan $b \in GF(2^8)$ \quad (dalam notasi Hex, $a = \text{'0D'}$ dan $b = \text{'06'}$)\

(i) $a + b = \text{'0D'} + \text{'06'} = (x^3 + x^2 + 1) + (x^2 + x) = x^3 + 2x^2 + x + 1 \pmod{2}$

\begin{align*}
\text{\color{red}Bagi tiap koefisien dengan 2, lalu ambil sisanya} \\
&= x^3 + 0x^2 + x + 1 \\
&= x^3 + x + 1 = 00001011 = \text{'0B'}
\end{align*}

Dalam representasi biner:

\begin{align*}
&00001101 \\
&00000110 + \\
&00001011 \rightarrow \text{sama dengan hasil operasi XOR}
\end{align*}

$\therefore a + b = 00001011 = a \text{ XOR } b$

\vspace{10mm}

\centering Rinaldi Munir/II4021 Kriptografi

\end{document}
